\documentclass{article}
\usepackage{amsmath,amssymb,amsthm}
\newcommand{\Om}{\Omega}
\newtheorem{lemma}{Lemma}
\begin{document}
\section{Decay for the heat equation}
Let $\Om\subset\mathbb R^N$ be a bounded smooth domain and let $\mu>0$
satisfy
\begin{equation}\label{eq:poincare-input}
 \int_{\Om}|\nabla v|^2\,dx\ge\mu\int_{\Om}|v|^2\,dx
 \qquad\text{for every }v\in H_0^1(\Om).
\end{equation}
This inequality is an assumption; no optimality of $\mu$ is needed.
\begin{lemma}\label{lem:heat-decay}
Let $T>0$ and let $u$ be a smooth real-valued solution of
$u_t=\Delta u$ on $[0,T]\times\Om$, with $u=0$ on
$[0,T]\times\partial\Om$. Then
\begin{equation}\label{eq:heat-conclusion}
 \|u(t)\|_{L^2(\Om)}\le e^{-\mu t}\|u(0)\|_{L^2(\Om)}
 \qquad(0\le t\le T).
\end{equation}
\end{lemma}
\begin{proof}
The squared $L^2$ norm measures the size of the solution. Define
\[
 A(t)=\int_{\Om}u(t,x)^2\,dx.
\]
Smoothness permits differentiation under the integral. Using the equation
and integrating by parts, we obtain
\[
 \begin{aligned}
 A'(t)
 &=2\int_{\Om}u(t,x)u_t(t,x)\,dx
 =2\int_{\Om}u(t,x)\Delta u(t,x)\,dx\\
 &=-2\int_{\Om}|\nabla u(t,x)|^2\,dx
 \le -2\mu A(t).
 \end{aligned}
\]
The boundary term vanishes because $u=0$ on $\partial\Om$, and the last
inequality follows from \eqref{eq:poincare-input} applied to $u(t,\cdot)$.
Consequently,
\[
 \frac{d}{dt}\bigl(e^{2\mu t}A(t)\bigr)
 =e^{2\mu t}\bigl(A'(t)+2\mu A(t)\bigr)\le0.
\]
Integration from $0$ to $t$ gives $A(t)\le e^{-2\mu t}A(0)$.
Taking nonnegative square roots yields \eqref{eq:heat-conclusion},
including times at which $A(t)=0$.
\end{proof}
\end{document}
